% 第十节：转录自“9.12中午_误差分析完成.docx”，优缺点分别采用“1、2、……”编号。
\subsection{模型优点}

\begin{enumerate}[label=\arabic*、,leftmargin=2em,labelsep=0pt,itemsep=0.5em,topsep=0.3em]
\item \textbf{模型具有较强的物理解释性。}
本文从药材内部热质传递机理出发，
逐步建立径向热质传递模型、变物性热质耦合模型以及收缩动边界模型，
能够较全面地描述烘干过程中温度、水分及尺寸变化规律。

\item \textbf{模型具有较好的递进性。}
针对不同阶段的研究需求，逐步引入物性变化、烘干判据和尺寸收缩因素，
在保证计算可行性的同时提高了模型与实际干燥过程的符合程度。

\item \textbf{数值求解过程稳定可靠。}
采用表面加密有限体积法和中点物性迭代处理模型计算，
并通过多种检验方法验证结果的稳定性，提高了模型预测精度。
\end{enumerate}

\subsection{模型缺点}

\begin{enumerate}[label=\arabic*、,leftmargin=2em,labelsep=0pt,itemsep=0.5em,topsep=0.3em]
\item \textbf{部分复杂物理过程进行了简化。}
模型中忽略了蒸发潜热、辐射换热等因素，
并采用一维径向假设描述药材内部传递过程，
对实际复杂结构药材的适用性存在一定限制。

\item \textbf{部分参数依赖题目给定或等效处理。}
水分扩散系数、换热系数等参数仍需要更精确的实验数据支持，
进一步完善参数测定可提高模型预测能力。
\end{enumerate}
